\documentclass[12pt,a4paper]{article}
\usepackage[utf8]{inputenc}
\usepackage[T5]{fontenc}
\usepackage[vietnamese]{babel}
\usepackage{amsmath,amssymb,amsfonts,amsthm,mathtools}
\usepackage{geometry}
\usepackage{enumitem}
\usepackage{array}
\usepackage{bm}
\usepackage{setspace}
\usepackage{titlesec}
\geometry{margin=2.2cm}
\setstretch{1.08}
\newcommand{\R}{\mathbb{R}}

\newcommand{\N}{\mathbb{N}}
\newcommand{\Eu}{\mathrm{Eu}}
\newcommand{\ip}[2]{\left\langle #1,#2\right\rangle}
\newcommand{\norm}[1]{\left\lVert #1\right\rVert}
\newcommand{\dist}{\mathrm{distance}}
\newcommand{\vect}[1]{\begin{pmatrix}#1\end{pmatrix}}
\titleformat{\section}{\large\bfseries}{\thesection.}{0.5em}{}
\setlength{\parindent}{0pt}
\setlength{\parskip}{5pt}
\begin{document}

\begin{center}
{\large\underline{Buổi 08 -- Ngày 13-09-2026 -- môn Đại số tuyến tính -- lớp MA003.F32.LT.CNTT}}
\end{center}

\textbf{Lưu ý:}
Cho $A$ là ma trận vuông cấp $n$ trên trường số thực $\R$.

Giả sử $A$ có các trị riêng thực là
\[
\begin{cases}
c_1=\lambda_1=\cdots\\
c_2=\lambda_2=\cdots\\
\vdots\\
c_n=\lambda_n=\cdots
\end{cases}
\]
Nghĩa là ta viết được
\[
p_A(\lambda)=\det(A-\lambda I_n)=(\lambda-\lambda_1)(\lambda-\lambda_2)\cdots(\lambda-\lambda_n).
\]
Từ đây ta sẽ tìm hệ số không chứa $\lambda$ để suy ra định thức của $A$.
Đây chính là tích của các trị riêng; nghĩa là
\[
\det(A)=\lambda_1\lambda_2\cdots\lambda_n.
\]

\underline{Ví dụ mẫu:} Cho $A$ là ma trận vuông thực, cấp $2$, có 2 trị riêng là $\lambda_1=-5,\lambda_2=4$. Tìm $\det(A)=?$\\
\underline{Giải:}

Ta có đa thức đặc trưng:
\[
\begin{aligned}
p_A(\lambda)&=\det(A-\lambda I_2)\\
&=\left|\begin{matrix}a_{11}&a_{12}\\a_{21}&a_{22}\end{matrix}
-\lambda\begin{matrix}1&0\\0&1\end{matrix}\right|
=(\lambda+5)(\lambda-4).
\end{aligned}
\]
Do đó
\[
\left|\begin{matrix}a_{11}-\lambda&a_{12}\\a_{21}&a_{22}-\lambda\end{matrix}\right|
=(\lambda+5)(\lambda-4),
\]
\[
(a_{11}-\lambda)(a_{22}-\lambda)-a_{12}a_{21}=\lambda^2+\lambda-20,
\]
\[
a_{11}a_{22}-\lambda a_{11}-\lambda a_{22}+\lambda^2-a_{12}a_{21}
=\lambda^2+\lambda-20,
\]
\[
a_{11}a_{22}-\lambda(a_{11}+a_{22})-a_{12}a_{21}=\lambda-20,
\]
\[
\begin{cases}
a_{11}+a_{22}=-1,\\
a_{11}a_{22}-a_{12}a_{21}=-20.
\end{cases}
\]
Suy ra
\[
\det(A)=a_{11}a_{22}-a_{12}a_{21}=-20.
\]
Kiểm chứng: $\det(A)=\lambda_1\lambda_2=(-5)\cdot4=-20$.

\section*{CHƯƠNG 5: KHÔNG GIAN EUCLIDE}
\subsection*{1/ MỘT SỐ KHÁI NIỆM}
Cho $V$ là không gian véc tơ (không gian tuyến tính) tổng quát trên trường số thực $\R$.

Một tích vô hướng (inner product -- cross product -- dot product) trên không gian $V$ là một ánh xạ
\[
\varphi:V\times V\to\R,
\qquad
(\alpha,\beta)\mapsto\varphi(\alpha,\beta)=\ip{\alpha}{\beta},
\quad \forall\alpha,\beta\in V,
\]
thỏa các tính chất:
\begin{enumerate}[label=\roman*)]
\item $\ip{\alpha}{\alpha}\ge0,\ \forall\alpha\in V$.

Ta có $\ip{\alpha}{\alpha}=0\Longleftrightarrow\alpha=0$.
\item $\ip{\beta}{\alpha}=\ip{\alpha}{\beta},\ \forall\alpha,\beta\in V$ (tính chất đối xứng).
\item $\ip{c\alpha_1+\alpha_2}{\beta}=c\ip{\alpha_1}{\beta}+\ip{\alpha_2}{\beta},\quad \forall c\in\R,\ \forall\alpha_1,\alpha_2,\beta\in V$.
\item $\ip{\alpha}{c\beta_1+\beta_2}=c\ip{\alpha}{\beta_1}+\ip{\alpha}{\beta_2},\quad \forall c\in\R,\ \forall\alpha,\beta_1,\beta_2\in V$.
\item $\ip{\alpha}{0}=\ip{0}{\alpha}=0,\ \forall\alpha\in V$.
\end{enumerate}

\textbf{Lưu ý:} Khi kiểm tra ánh xạ có phải là một tích vô hướng trên không gian $V$ hay không, thì ta chỉ cần kiểm 3 tính chất i), ii), iii), trong đó tính chất i) là khó nhất. Tính chất iv) mặc nhiên thỏa nếu ta có tính chất ii) và iii); còn tính chất v) là luôn thỏa với mọi véc tơ nên ta không cần kiểm tính chất v).

Để chứng minh ánh xạ không phải là một tích vô hướng trên không gian thì ta chỉ cần chứng tỏ 1 trong 4 tính chất đầu tiên (i), ii), iii), iv) nêu trên bị vi phạm.

Ta gọi không gian véc tơ $V$, có trang bị thêm một tích vô hướng đi kèm là không gian Euclide (Euclidian space), và ký hiệu là
\[
\Eu=(V,\ip{\cdot}{\cdot}).
\]

\underline{Ví dụ mẫu 1:} Trên $V=\R^3$, cho ánh xạ
\[
\varphi:\R^3\times\R^3\to\R
\]
với
\[
\alpha=\vect{x_1\\x_2\\x_3},\quad
\beta=\vect{y_1\\y_2\\y_3},
\]
\[
\ip{\alpha}{\beta}=2x_1y_1-x_1y_2+3x_3y_3-x_2y_1+4x_2y_2.
\]
Hỏi $\ip{\alpha}{\beta}$ có phải là tích vô hướng hay không? Vì sao?

\underline{Ví dụ mẫu 2:} Cho $V=\R^2$, cho ánh xạ
\[
\varphi:\R^2\times\R^2\to\R
\]
với
\[
\alpha=\vect{x_1\\x_2},\quad \beta=\vect{y_1\\y_2},
\]
\[
\ip{\alpha}{\beta}=x_1y_1-5x_1y_2+2x_2y_1-10x_2y_2+3y_1y_2+4x_1^2.
\]
Hỏi $\ip{\alpha}{\beta}$ có phải là tích vô hướng hay không? Vì sao?

\underline{Giải:}

\underline{Ví dụ mẫu 1:} Ta chứng tỏ $\ip{\alpha}{\beta}$ là một tích vô hướng như sau.

\textit{i/} Ta có
\[
\begin{aligned}
\ip{\alpha}{\alpha}
&=2x_1^2-x_1x_2+3x_3^2-x_2x_1+4x_2^2\\
&=2x_1^2-2x_1x_2+4x_2^2+3x_3^2\\
&=(x_1^2-2x_1x_2+x_2^2)+x_1^2+3x_2^2+3x_3^2\\
&=(x_1-x_2)^2+x_1^2+3x_2^2+3x_3^2\ge0.
\end{aligned}
\]
là luôn đúng. Nên ta có $\ip{\alpha}{\alpha}\ge0,\ \forall\alpha\in\R^3$. (*)

\textit{ii/} Ta có
\[
\begin{aligned}
\ip{\alpha}{\beta}
&=2x_1y_1-x_1y_2+3x_3y_3-x_2y_1+4x_2y_2,
\\
\ip{\beta}{\alpha}
&=2y_1x_1-y_1x_2+3y_3x_3-y_2x_1+4y_2x_2\\
&=2x_1y_1-x_2y_1+3x_3y_3-x_1y_2+4x_2y_2\\
&=2x_1y_1-x_1y_2+3x_3y_3-x_2y_1+4x_2y_2.
\end{aligned}
\]
Suy ra $\ip{\alpha}{\beta}=\ip{\beta}{\alpha},\ \forall\alpha,\beta\in\R^3$. (**)

\textit{iii/} Ta chứng tỏ tuyến tính theo biến thứ nhất, nghĩa là cần chứng minh
\[
\ip{c\alpha_1+\alpha_2}{\beta}=c\ip{\alpha_1}{\beta}+\ip{\alpha_2}{\beta}.
\]
Với
\[
\alpha_1=\vect{x_1\\x_2\\x_3},\quad
\alpha_2=\vect{y_1\\y_2\\y_3},\quad
\beta=\vect{z_1\\z_2\\z_3},
\]
ta có
\[
c\alpha_1+\alpha_2=\vect{cx_1+y_1\\cx_2+y_2\\cx_3+y_3}.
\]
Do đó
\[
\begin{aligned}
\ip{c\alpha_1+\alpha_2}{\beta}
&=2(cx_1+y_1)z_1-(cx_1+y_1)z_2+3(cx_3+y_3)z_3-(cx_2+y_2)z_1+4(cx_2+y_2)z_2\\
&=c(2x_1z_1-x_1z_2+3x_3z_3-x_2z_1+4x_2z_2)\\
&\quad +(2y_1z_1-y_1z_2+3y_3z_3-y_2z_1+4y_2z_2)\\
&=c\ip{\alpha_1}{\beta}+\ip{\alpha_2}{\beta}.
\end{aligned}
\]
(***). Từ (*), (**), (***), suy ra $\ip{\alpha}{\beta}$ là một tích vô hướng trên $\R^3$.

\underline{Ví dụ mẫu 2:} Ta chứng tỏ $\ip{\alpha}{\beta}$ không phải là một tích vô hướng.

\textit{Cách 1:} Ta chọn $\alpha=(0,4)^T$ thì
\[
\ip{\alpha}{\alpha}=0-5\cdot4+2\cdot4-10\cdot4+3\cdot4+4\cdot4^2=-160<0.
\]
Suy ra không phải là một tích vô hướng trên $\R^2$.

\textit{Cách 2:} Ta chọn
\[
\alpha=\vect{2\\-1},\qquad \beta=\vect{-3\\4}.
\]
Suy ra
\[
\ip{\alpha}{\beta}=2(-3)-5\cdot2\cdot4+2(-1)(-3)-10(-1)4+3(-3)4+4\cdot2^2=-20,
\]
Mà
\[
\ip{\beta}{\alpha}=(-3)2-5(-3)(-1)+2\cdot4\cdot2-10\cdot4(-1)+3\cdot2(-1)+4(-3)^2=65.
\]
Suy ra $\ip{\alpha}{\beta}\ne\ip{\beta}{\alpha}$ cho nên không phải là một tích vô hướng trên $\R^2$.

\textit{Cách 3:} Ta chọn
\[
c=2,\quad \alpha_1=\vect{-1\\3},\quad \alpha_2=\vect{2\\1},\quad \beta=\vect{-4\\5}.
\]
Suy ra
\[
c\alpha_1+\alpha_2=2\vect{-1\\3}+\vect{2\\1}=\vect{0\\7}.
\]
\[
\ip{c\alpha_1+\alpha_2}{\beta}=-466,
\]
trong khi
\[
c\ip{\alpha_1}{\beta}=-402,\qquad \ip{\alpha_2}{\beta}=-160.
\]
Do đó
\[
c\ip{\alpha_1}{\beta}+\ip{\alpha_2}{\beta}=-562\ne-466=\ip{c\alpha_1+\alpha_2}{\beta},
\]
cho nên không phải là một tích vô hướng trên $\R^2$.

\textit{Cách 4:} Ta chọn
\[
c=-1,\quad \alpha=\vect{2\\3},\quad \beta_1=\vect{-2\\5},\quad \beta_2=\vect{1\\-6}.
\]
Suy ra
\[
c\beta_1+\beta_2=-\vect{-2\\5}+\vect{1\\-6}=\vect{3\\-11}.
\]
\[
\ip{\alpha}{c\beta_1+\beta_2}=381,
\]
trong khi
\[
c\ip{\alpha}{\beta_1}=230,\qquad \ip{\alpha}{\beta_2}=246.
\]
Vì $230+246=476\ne381$, nên tính tuyến tính theo biến thứ hai bị vi phạm; không phải là một tích vô hướng trên $\R^2$.

\underline{Bài tập tương tự:} Kiểm tra xem ánh xạ có phải là một tích vô hướng trên không gian $V$ hay không? Vì sao?
\begin{enumerate}[label=\alph*/]
\item $V=\R^2$ và $\ip{\alpha}{\beta}=5x_1y_1+8x_2y_2$.
\item $V=\R^2$ và $\ip{\alpha}{\beta}=6x_1y_2-9x_2y_1+2x_1^2$.
\item $V=\R^3$ và
\[
\ip{\alpha}{\beta}=x_1y_1+2x_1y_2+3x_3y_3+2x_2y_1+12x_2y_2.
\]
\item $V=\R^3$ và
\[
\ip{\alpha}{\beta}=6x_1y_1-5x_2y_3+2x_2y_1-8x_2y_2.
\]
\item $V=\R^3$ và
\[
\ip{\alpha}{\beta}=4x_1y_1-2x_2y_3+3x_2y_2-7x_1x_3.
\]
\end{enumerate}

\textbf{Lưu ý:} Trên không gian $V=\R^n$, ta luôn có tích vô hướng tiêu chuẩn:
\[
\ip{\alpha}{\beta}=x_1y_1+x_2y_2+\cdots+x_ny_n.
\]

\subsection*{2/ ĐỘ DÀI VÉC TƠ VÀ KHOẢNG CÁCH GIỮA 2 VÉC TƠ}
Cho không gian Euclide $\Eu=(V,\ip{\cdot}{\cdot})$. Cho 2 véc tơ $\alpha,\beta\in\Eu$.
\[
\norm{\alpha}=\sqrt{\ip{\alpha}{\alpha}}=\text{độ dài (module) của véc tơ }\alpha.
\]
Và khoảng cách giữa 2 véc tơ $\alpha,\beta$ là
\[
d(\alpha,\beta)=\norm{\alpha-\beta}=\sqrt{\ip{\alpha-\beta}{\alpha-\beta}}=\dist(\alpha,\beta)
\]
\[
=\norm{\beta-\alpha}=\sqrt{\ip{\beta-\alpha}{\beta-\alpha}}.
\]

\underline{Ví dụ 1:} Trên không gian Euclide $\Eu=(V=\R^4,\ip{\cdot}{\cdot})$ với $\ip{\cdot}{\cdot}$ là tích vô hướng tiêu chuẩn trên $\R^4$, cho các véc tơ
\[
\alpha_1=(3,-2,1,5),\quad \alpha_2=(-1,7,3,-8),\quad \alpha_3=(6,-4,2,-1).
\]
Tìm độ dài các véc tơ và tính khoảng cách giữa các véc tơ; nghĩa là tìm
\[
\left\{\begin{array}{l}
\norm{\alpha_1},\norm{\alpha_2},\norm{\alpha_3},\\
d(\alpha_1,\alpha_2),d(\alpha_2,\alpha_3),d(\alpha_1,\alpha_3).
\end{array}\right.
\]
Gợi ý giải:
\[
\norm{\alpha_1}=\sqrt{\ip{\alpha_1}{\alpha_1}}=\sqrt{3^2+(-2)^2+1^2+5^2}=\sqrt{39}.
\]
Và
\[
\alpha_1-\alpha_2=(4,-9,-2,13),
\]
\[
d(\alpha_1,\alpha_2)=\sqrt{4^2+(-9)^2+(-2)^2+13^2}=3\sqrt{30}.
\]

\underline{Ví dụ 2:} Trên không gian Euclide $\Eu=(V=\R^5,\ip{\cdot}{\cdot})$ với tích vô hướng tiêu chuẩn trên $\R^5$, cho
\[
\alpha_1=(4,1,-6,7,5),\quad \alpha_2=(2,-3,-1,4,-5),\quad \alpha_3=(-7,2,-2,9,3).
\]
Tìm độ dài các véc tơ và tính khoảng cách giữa các véc tơ; nghĩa là tìm
\[
\left\{\begin{array}{l}
\norm{\alpha_1},\norm{\alpha_2},\norm{\alpha_3},\\
d(\alpha_1,\alpha_2),d(\alpha_2,\alpha_3),d(\alpha_1,\alpha_3).
\end{array}\right.
\]

\subsection*{3/ TRỰC GIAO (ORTHOGONAL) VÀ TRỰC CHUẨN (ORTHONORMAL)}
Cho không gian Euclide $\Eu=(V,\ip{\cdot}{\cdot})$.
Cho tập hợp $S=\{\alpha_1,\alpha_2,\ldots,\alpha_m\}$ trên $\Eu$.

Ta gọi 2 véc tơ $\alpha_1$ và $\alpha_2$ là vuông góc nhau, và kí hiệu $\alpha_1\perp\alpha_2$ khi và chỉ khi
\[
\ip{\alpha_1}{\alpha_2}=0.
\]
Nếu các véc tơ trong $S$ đều vuông góc với nhau, nghĩa là
\[
\ip{\alpha_i}{\alpha_j}=0,\quad\forall\alpha_i,\alpha_j\in S,\ i\ne j,
\]
thì ta nói $S$ là tập hợp trực giao (orthogonal set).

Nếu $S$ là một tập hợp trực giao và có thêm tính chất là độ dài các véc tơ trong $S$ đều bằng $1$ thì ta nói $S$ là tập hợp trực chuẩn (orthonormal set).

Cho véc tơ $\alpha\in\Eu$. Ta nói $\alpha\perp S$ nếu
\[
\ip{\alpha}{\alpha_i}=0,\quad\forall\alpha_i\in S.
\]
Ta gọi tập hợp chứa các véc tơ vuông góc với tập hợp $S$ là không gian trực giao với $S$ và kí hiệu
\[
S^\perp=\{\alpha\in\Eu\mid \alpha\perp S\}.
\]

\subsection*{4/ TRỰC GIAO HÓA VÀ TRỰC CHUẨN HÓA MỘT TẬP HỢP BẰNG PHƯƠNG PHÁP GRAM--SCHMIDT}
Cho không gian Euclide $\Eu=(V,\ip{\cdot}{\cdot})$. Cho tập hợp $S=\{\alpha_1,\alpha_2,\ldots,\alpha_m\}$ trên $\Eu$.
Ta trực giao hóa $S$ thành tập hợp trực giao $S'$ trên $\Eu$ như sau:
\[
\beta_1=\alpha_1,
\]
\[
\beta_2=\alpha_2-\frac{\ip{\alpha_2}{\beta_1}}{\norm{\beta_1}^2}\beta_1,
\]
\[
\beta_3=\alpha_3-\frac{\ip{\alpha_3}{\beta_2}}{\norm{\beta_2}^2}\beta_2
-\frac{\ip{\alpha_3}{\beta_1}}{\norm{\beta_1}^2}\beta_1,
\]
\[
\beta_4=\alpha_4-\frac{\ip{\alpha_4}{\beta_3}}{\norm{\beta_3}^2}\beta_3
-\frac{\ip{\alpha_4}{\beta_2}}{\norm{\beta_2}^2}\beta_2
-\frac{\ip{\alpha_4}{\beta_1}}{\norm{\beta_1}^2}\beta_1,
\]
\[
\vdots
\]
\[
\beta_m=\alpha_m-\frac{\ip{\alpha_m}{\beta_{m-1}}}{\norm{\beta_{m-1}}^2}\beta_{m-1}-\cdots
-\frac{\ip{\alpha_m}{\beta_2}}{\norm{\beta_2}^2}\beta_2
-\frac{\ip{\alpha_m}{\beta_1}}{\norm{\beta_1}^2}\beta_1.
\]
Suy ra $S'=\{\beta_1,\beta_2,\ldots,\beta_m\}$ là một tập hợp trực giao cần tìm.

Ta trực chuẩn hóa tiếp theo thành tập hợp trực chuẩn $S''$ như sau:
\[
\gamma_1=\frac{\beta_1}{\norm{\beta_1}},\qquad
\gamma_2=\frac{\beta_2}{\norm{\beta_2}},\qquad
\ldots,\qquad
\gamma_m=\frac{\beta_m}{\norm{\beta_m}}.
\]
Suy ra $S''=\{\gamma_1,\gamma_2,\ldots,\gamma_m\}$ là một tập hợp trực chuẩn trên $\Eu$.

\underline{Ví dụ mẫu 3:} Trên không gian Euclide $\Eu=(V=\R^3,\ip{\cdot}{\cdot})$, với tích vô hướng tiêu chuẩn trên $\R^3$, cho tập hợp
\[
S=\{u_1=(1,1,1),u_2=(-1,1,0),u_3=(1,2,1)\}.
\]
Hãy trực chuẩn hóa $S$.

\underline{Giải:} Dùng pp Gram--Schmidt, ta làm như sau:
\[
v_1=u_1=(1,1,1),
\]
\[
\begin{aligned}
v_2&=u_2-\frac{\ip{u_2}{v_1}}{\norm{v_1}^2}v_1\\
&=(-1,1,0)-\frac{(-1)(1)+1\cdot1+0\cdot1}{1^2+1^2+1^2}(1,1,1)\\
&=(-1,1,0)-0(1,1,1)=(-1,1,0),
\end{aligned}
\]
\[
\begin{aligned}
v_3&=u_3-\frac{\ip{u_3}{v_2}}{\norm{v_2}^2}v_2-\frac{\ip{u_3}{v_1}}{\norm{v_1}^2}v_1\\
&=(1,2,1)-\frac{1(-1)+2(1)+1(0)}{(-1)^2+1^2+0^2}(-1,1,0)\\
&\quad-\frac{1\cdot1+2\cdot1+1\cdot1}{1^2+1^2+1^2}(1,1,1)\\
&=(1,2,1)-\frac12(-1,1,0)-\frac43(1,1,1)\\
&=\left(\frac16,\frac16,-\frac13\right).
\end{aligned}
\]
Ta có $S'=\{v_1,v_2,v_3\}$ là một tập hợp trực giao trên $\R^3$.
Đặt
\[
q_1=\frac{v_1}{\norm{v_1}}=\frac{(1,1,1)}{\sqrt3}=\left(\frac1{\sqrt3},\frac1{\sqrt3},\frac1{\sqrt3}\right),
\]
\[
q_2=\frac{v_2}{\norm{v_2}}=\frac{(-1,1,0)}{\sqrt2}=\left(-\frac1{\sqrt2},\frac1{\sqrt2},0\right),
\]
\[
\begin{aligned}
q_3&=\frac{v_3}{\norm{v_3}}
=\frac{(1/6,1/6,-1/3)}{\sqrt{(1/6)^2+(1/6)^2+(-1/3)^2}}\\
&=\frac1{\sqrt6}\left(\frac16,\frac16,-\frac13\right)\cdot 6
=\left(\frac1{\sqrt6},\frac1{\sqrt6},-\frac2{\sqrt6}\right).
\end{aligned}
\]
Suy ra $S''=\{q_1,q_2,q_3\}$ là một tập hợp trực chuẩn của $\R^3$.

\underline{Bài tập tương tự:} Yêu cầu như ví dụ mẫu 3. Trên không gian Euclide $\Eu=(V=\R^3,\ip{\cdot}{\cdot})$, với tích vô hướng tiêu chuẩn trên $\R^3$, cho tập hợp
\[
S=\{u_1=(2,3,6),u_2=(5,-3,8),u_3=(8,5,3)\}.
\]
Hãy trực chuẩn hóa $S$.

\underline{Bài tập tương tự khác:} Trên không gian Euclide $\Eu=(V=\R^3,\ip{\cdot}{\cdot})$ với
\[
\ip{\alpha}{\beta}=x_1y_1+3x_2y_2+2x_3y_3,
\]
là một tích vô hướng trên $\R^3$. Hãy trực chuẩn hóa tập hợp
\[
S=\{u_1=(1;0;1),u_2=(-4;0;2),u_3=(-3;-1;-3)\}.
\]

\underline{Bài tập tương tự:} hãy trực chuẩn hóa các tập hợp sau trên không gian Euclide tương ứng.
\begin{enumerate}
\item $\Eu=(\R^3,\ip{\cdot}{\cdot})$ và $S=\{\alpha_1=(2;3;6),\alpha_2=(-3;6;-2),\alpha_3=(6;2;-3)\}$.
\item $\Eu=(\R^3,\ip{\cdot}{\cdot})$ và $S=\{\alpha_1=(1;2;2),\alpha_2=(-2;2;-1),\alpha_3=(2;1;-2)\}$.
\item $\Eu=(\R^3,\ip{\cdot}{\cdot})$ và $S=\{\alpha_1=(8;1;-4),\alpha_2=(1;0;2),\alpha_3=(-2;8;-2)\}$.
\end{enumerate}

\subsection*{5/ TỌA ĐỘ VÉC TƠ THEO CƠ SỞ TRỰC CHUẨN}
Cho không gian Euclide $\Eu=(V_n,\ip{\cdot}{\cdot})$ có cơ sở trực chuẩn $\Xi=\{\gamma_1,\gamma_2,\ldots,\gamma_n\}$.

Cho trước véc tơ $\alpha\in V_n$, thì ta luôn có $c_1,c_2,\ldots,c_n$ sao cho
\[
\alpha=c_1\gamma_1+c_2\gamma_2+\cdots+c_n\gamma_n.
\]
Trong đó
\[
c_1=\ip{\alpha}{\gamma_1},\quad c_2=\ip{\alpha}{\gamma_2},\quad\ldots,\quad c_n=\ip{\alpha}{\gamma_n}.
\]
Suy ra, tọa độ của véc tơ $\alpha$ theo cơ sở trực chuẩn $\Xi$ là
\[
[\alpha]_{\Xi}=\begin{pmatrix}
\ip{\alpha}{\gamma_1}\\
\ip{\alpha}{\gamma_2}\\
\vdots\\
\ip{\alpha}{\gamma_n}
\end{pmatrix}.
\]

\underline{Ví dụ mẫu:} Trên không gian Euclide $\Eu=(V=\R^3,\ip{\cdot}{\cdot})$ có cơ sở trực chuẩn
\[
\Xi=\left\{\gamma_1=\frac13(1,2,2),\gamma_2=\frac13(-2,2,-1),\gamma_3=\frac13(2,1,-2)\right\}.
\]
a/ Cho $\alpha=(6;-7;1)\in\R^3$. Tìm $[\alpha]_{\Xi}=?$\\
b/ Cho $\beta=(-2;1;-3)\in\R^3$. Tìm $[\beta]_{\Xi}=?$

Giải: Do $\Xi$ là cơ sở trực chuẩn của $\Eu=(\R^3,\ip{\cdot}{\cdot})$, nên ta có:
\[
c_1=\ip{\alpha}{\gamma_1}=6\cdot\frac13+(-7)\frac23+1\frac23=-2,
\]
\[
c_2=\ip{\alpha}{\gamma_2}=6\left(-\frac23\right)+(-7)\frac23+1\left(-\frac13\right)=-9,
\]
\[
c_3=\ip{\alpha}{\gamma_3}=6\frac23+(-7)\frac13+1\left(-\frac23\right)=1.
\]
Suy ra
\[
[\alpha]_{\Xi}=\begin{pmatrix}-2\\-9\\1\end{pmatrix}.
\]
Với $\beta$:
\[
c_1=\ip{\beta}{\gamma_1}=(-2)\frac13+1\frac23+(-3)\frac23=-2,
\]
\[
c_2=\ip{\beta}{\gamma_2}=(-2)\left(-\frac23\right)+1\frac23+(-3)\left(-\frac13\right)=3,
\]
\[
c_3=\ip{\beta}{\gamma_3}=(-2)\frac23+1\frac13+(-3)\left(-\frac23\right)=1.
\]
Suy ra
\[
[\beta]_{\Xi}=\begin{pmatrix}-2\\3\\1\end{pmatrix}.
\]

\subsection*{6/ CHÉO HÓA TRỰC GIAO MA TRẬN VUÔNG (phần đọc thêm cho vui)}
Cho ma trận thực, vuông $A$, cấp $n$.
Ta gọi $A$ là ma trận trực giao, nếu và chỉ nếu $A$ khả nghịch và
\[
A^{-1}=A^T.
\]
Ta chéo hóa trực giao ma trận $A$ như sau:
\begin{itemize}
\item Ta tìm
\[
p_A(x)=p_A(\lambda)=\det(xI_n-A)=\det(A-\lambda I_n).
\]
\item Giải pt $p_A(x)=0$ hoặc $p_A(\lambda)=0$ trên $\R$ để tìm nghiệm.
\item Nếu pt vô nghiệm thì ta kết luận không thể chéo hóa $A$ trên $\R$ và ta dừng bài toán.
\item Nếu pt có các nghiệm $x_1=\lambda_1=\cdots,\ x_2=\lambda_2=\cdots,\ldots,x_k=\lambda_k=\cdots$ thì ta nói $p_A(x)$ tách được trên $\R$ và viết
\[
p_A(x)=(x-c_1)^{r_1}(x-c_2)^{r_2}(x-c_3)^{r_3}\cdots(x-c_k)^{r_k},
\]
trong đó $c_1=\lambda_1,\ldots,c_k=\lambda_k$.
\[
r_1=\cdots,\quad r_2=\cdots,\quad\ldots,\quad r_k=\cdots.
\]
Suy ra nếu $r_1=r_2=\cdots=r_k=1$ thì $A$ chéo hóa được trên $\R$.
\item Tiếp theo, ta tìm cơ sở $\alpha_j$ cho các không gian riêng
\[
E_{c_j}=E_{\lambda_j}=\{X\in\R^n\mid(A-c_jI_n)X=0\}
\]
và xác định số chiều $\dim_{\R}E_{c_j}=?$.
\end{itemize}
Nếu có (ít nhất) 1 số $j\in\{1,2,\ldots,k\}$ mà $\dim_{\R}E_{c_j}<r_j$ thì $A$ không chéo hóa được và ta dừng bài toán.
Nếu $\dim_{\R}E_{c_j}=r_j$ với mọi $j=1,2,\ldots,k$ thì $A$ chéo hóa được trên $\R$.

Ta trực chuẩn hóa cơ sở $\alpha^j$ bằng pp Gram--Schmidt thành cơ sở trực chuẩn $\Xi_j$.
Gọi $\Xi=\Xi_1\cup\Xi_2\cup\cdots\cup\Xi_k$ thì $\Xi$ là một cơ sở trực chuẩn của $\R^n$.
Đặt $P=I=P_{\beta_0\to\Xi}$ thì $P=I$ là ma trận trực giao, thỏa
\[
D=P^{-1}AP=P^TAP=\begin{pmatrix}
c_1&&\\&c_2&\\&&\ddots\\&&&c_n
\end{pmatrix}.
\]

\underline{Ví dụ mẫu:} Cho ma trận vuông thực
\[
A=\begin{pmatrix}
2&-1&-1\\
-1&2&-1\\
-1&-1&2
\end{pmatrix}.
\]
Hãy chéo hóa trực giao ma trận $A$.

\underline{Giải:} Ta có
\[
p_A(x)=\det(xI_3-A)=
\left|\begin{matrix}x-2&1&1\\1&x-2&1\\1&1&x-2\end{matrix}\right|
=(x-2)^3+1+1-(x-2)-(x-2)-(x-2)=x(x-3)^2.
\]
Suy ra
\[
c_1=\lambda_1=0,\ r_1=1,\qquad c_2=\lambda_2=3,\ r_2=2.
\]

Tiếp theo, ta tìm cơ sở cho các không gian riêng.

Với $c_1=\lambda_1=0$, ta có
\[
E_0=\{X\in\R^3\mid(A-0I_3)X=0\}.
\]
Giải hệ pt $(A-0I_3)X=0$. Phép biến đổi sơ cấp cho dạng bậc thang
\[
\left[\begin{array}{ccc|c}
2&-1&-1&0\\
-1&2&-1&0\\
-1&-1&2&0
\end{array}\right]
\longrightarrow
\left[\begin{array}{ccc|c}
2&-1&-1&0\\
0&3&-3&0\\
0&0&0&0
\end{array}\right].
\]
Do cột 3 không bán chuẩn hóa được nên hệ có vô số nghiệm. Đặt $x_3=t$, $t\in\R$.
\[
\begin{cases}
2x_1-x_2-x_3=0,\\
3x_2-3x_3=0
\end{cases}
\Longrightarrow
\begin{cases}
2x_1=x_2+x_3=2t,\\x_2=x_3=t.
\end{cases}
\]
Nên ta có nghiệm $x_1=x_2=x_3=t$, $\forall t\in\R$.
Suy ra cơ sở $\alpha_1=\{(1;1;1)\}$ và số chiều là $\dim_{\R}E_{c_1}=1=r_1$.

Với $c_2=\lambda_2=3$, ta có
\[
E_3=\{X\in\R^3\mid(A-3I_3)X=0\}.
\]
Giải hệ $(A-3I_3)X=0$ được
\[
\left[\begin{array}{ccc|c}
-1&-1&-1&0\\
-1&-1&-1&0\\
-1&-1&-1&0
\end{array}\right]
\longrightarrow
\left[\begin{array}{ccc|c}
-1&-1&-1&0\\
0&0&0&0\\
0&0&0&0
\end{array}\right].
\]
\[
x_1=-a-b,\qquad x_2=a,\qquad x_3=b,\quad \forall a,b\in\R.
\]
Suy ra cơ sở $\alpha_2=\{(-1;1;0),(-1;0;1)\}$ và số chiều là $\dim_{\R}E_{c_2}=2=r_2$.

Từ (1), (2) suy ra $A$ chéo hóa được trên $\R$.

Ta trực chuẩn hóa bằng pp Gram--Schmidt cho cơ sở $\alpha_1$ và $\alpha_2$ như sau.

$a_1$ thành
\[
\Xi_1=\left\{\gamma_1=\frac{\alpha_1}{\norm{\alpha_1}}=\frac{(1,1,1)}{\sqrt{1^2+1^2+1^2}}=\frac1{\sqrt3}(1,1,1)\right\}.
\]

$a_2$ thành
\[
\beta=\left\{\beta_2=\alpha_2=(-1,1,0),\quad
\beta_3=\alpha_3-\frac{\ip{\alpha_3}{\beta_2}}{\norm{\beta_2}^2}\beta_2\right\}
\]
\[
=\left\{\beta_2=(-1,1,0),\quad
\beta_3=(-1,0,1)-\frac{(-1)(-1)+0\cdot1+1\cdot0}{(-1)^2+1^2+0^2}(-1,1,0)\right\}
\]
\[
=\left\{\beta_2=(-1,1,0),\quad\beta_3=\frac12(-1,-1,2)\right\}.
\]
Trực chuẩn hóa $\beta$ thành
\[
\Xi_2=\left\{\gamma_2=\frac{\beta_2}{\norm{\beta_2}}=\frac1{\sqrt2}(-1,1,0),\quad
\gamma_3=\frac{\beta_3}{\norm{\beta_3}}=\frac1{\sqrt6}(-1,-1,2)\right\}.
\]
Đặt $\Xi=\Xi_1\cup\Xi_2=\{\gamma_1,\gamma_2,\gamma_3\}$ là một cơ sở trực chuẩn của $\R^3$.

Gọi
\[
P=T=P_{\beta_0\to\Xi}=
\begin{pmatrix}
1/\sqrt3&-1/\sqrt2&-1/\sqrt6\\
1/\sqrt3&1/\sqrt2&-1/\sqrt6\\
1/\sqrt3&0&2/\sqrt6
\end{pmatrix}.
\]
Thì $P=T$ là ma trận trực giao, nghĩa là $P^{-1}=P^T$, thỏa
\[
D=P^{-1}AP=P^TAP=
\begin{pmatrix}0&0&0\\0&3&0\\0&0&3\end{pmatrix}.
\]
Đây là dạng chéo hóa trực giao cần tìm của $A$.

\underline{Bài tập tương tự:} chéo hóa trực giao ma trận vuông thực
\[
A=\begin{pmatrix}1&0&1\\0&1&0\\1&0&1\end{pmatrix}.
\]

\section*{NỘI DUNG ÔN TẬP CUỐI KỲ -- THỜI GIAN LÀM BÀI 90 PHÚT}
\textbf{SV chỉ được sử dụng tài liệu GIẤY}

\underline{Câu 1.} Trên $V=\R^5$ (hoặc $V=\R^6$), cho tập hợp
\[
W=\{\cdots\mid\cdots\}.
\]
Chứng minh rằng $W$ là không gian véc tơ con của $V$. Tìm hệ sinh, cơ sở và xác định số chiều cho $W$.

\underline{Câu 2.} Trên $V=\R^3$, cho tập hợp
\[
a=\{u_1=\cdots;u_2=\cdots;u_3=\cdots\}.
\]
\begin{enumerate}[label=\alph*/]
\item Chứng minh rằng $a$ là một cơ sở của $\R^3$.
\item Trực chuẩn hóa cơ sở $a$ thành cơ sở trực chuẩn $b$.
\item Tìm tọa độ của một véc tơ $v=\cdots\in\R^3$ theo cơ sở $a$.
\item Tìm tọa độ của một véc tơ $v=\cdots\in\R^3$ theo cơ sở $b$.
\end{enumerate}

\underline{Câu 3:} Trên $V=\R^4$, với tích vô hướng $\ip{\cdot}{\cdot}=\cdots$, cho hệ
$S=\{u_1=\cdots;u_2=\cdots;u_3=\cdots\}$. Dùng pp Gram--Schmidt để trực chuẩn hóa $S$.
Cho véc tơ $\lambda=\cdots$. Tìm tọa độ của véc tơ $\lambda$ theo cơ sở vừa trực chuẩn hóa được.

\underline{Câu 4.} Cho ma trận vuông thực, cấp $2$, $A=\cdots$. Hãy chéo hóa ma trận $A$, rồi tìm $A^m$, $\forall m\in\N$ và suy ra $A^{2026}$.

\begin{center}
----------HẾT----------
\end{center}

\end{document}
